\documentclass[11pt,oneside]{book}

\usepackage[a4paper,margin=1.05in]{geometry}
\usepackage[T1]{fontenc}
\usepackage[utf8]{inputenc}
\usepackage{lmodern}
\usepackage{microtype}
\usepackage{amsmath,amssymb,amsthm,mathtools}
\usepackage{bm}
\usepackage{booktabs}
\usepackage{enumitem}
\usepackage{xcolor}
\usepackage{graphicx}
\usepackage{hyperref}
\usepackage{tikz}
\usepackage{pgfplots}
\usepackage{listings}
\usepackage{csquotes}
\usepackage{fancyhdr}

\usetikzlibrary{arrows.meta,calc,decorations.pathmorphing,positioning,shapes.geometric,fit,backgrounds}
\pgfplotsset{compat=1.18}

\definecolor{ink}{HTML}{1B1E28}
\definecolor{deepblue}{HTML}{224B8D}
\definecolor{teal}{HTML}{1A7F78}
\definecolor{amber}{HTML}{B36B00}
\definecolor{crimson}{HTML}{A4333A}
\definecolor{violet}{HTML}{5A4B9B}
\definecolor{softgray}{HTML}{F3F5F7}
\definecolor{linegray}{HTML}{C8CED8}

\hypersetup{
  colorlinks=true,
  linkcolor=deepblue,
  citecolor=teal,
  urlcolor=violet,
  pdftitle={Adiabatic Hypercomplex Learning Towers},
  pdfauthor={Jesús Vilela Jato and repository agents}
}

\lstdefinestyle{code}{
  basicstyle=\ttfamily\small,
  backgroundcolor=\color{softgray},
  frame=single,
  rulecolor=\color{linegray},
  breaklines=true,
  columns=fullflexible
}
\lstset{style=code}

\pagestyle{fancy}
\fancyhf{}
\lhead{\nouppercase{\leftmark}}
\rhead{\thepage}
\renewcommand{\headrulewidth}{0.4pt}

\theoremstyle{definition}
\newtheorem{definition}{Definition}[chapter]
\newtheorem{assumption}[definition]{Assumption}
\newtheorem{protocol}[definition]{Protocol}
\newtheorem{status}[definition]{Status Label}
\theoremstyle{plain}
\newtheorem{theorem}[definition]{Theorem}
\newtheorem{proposition}[definition]{Proposition}
\newtheorem{lemma}[definition]{Lemma}
\theoremstyle{remark}
\newtheorem{remark}[definition]{Remark}

\newcommand{\Sop}[1]{\mathsection\mathrm{#1}}
\newcommand{\Return}{\Sop{RETURN}}
\newcommand{\Ground}{\Sop{GROUND}}
\newcommand{\Lift}{\Sop{LIFT}}
\newcommand{\Immersion}{\Sop{IMMERSION}}
\newcommand{\Resonate}{\Sop{RESONATE}}
\newcommand{\Recognize}{\Sop{RECOGNIZE}}
\newcommand{\Telos}{\mathrm{TELOS}}
\newcommand{\Repo}{X_{\mathrm{repo}}}
\newcommand{\Usable}{S_{\mathrm{usable}}}

\title{\Huge\bfseries Adiabatic Hypercomplex Learning Towers\\[0.4em]
\Large A Tectonic Thesis on Immersion, Mutual Recognition, TELOS, and Grounded Return}
\author{Jesús Vilela Jato\\\small with repository-grounded agentic formalization}
\date{May 2026}

\begin{document}

\frontmatter
\begin{titlepage}
\centering
\vspace*{1.2cm}
{\Huge\bfseries Adiabatic Hypercomplex\\Learning Towers\par}
\vspace{0.5cm}
{\Large A Tectonic Thesis on Immersion, Mutual Recognition, TELOS, and Grounded Return\par}
\vspace{1.2cm}
\begin{tikzpicture}[scale=1.0]
  \fill[softgray] (0,0) circle (3.1);
  \draw[linegray, line width=0.8pt] (0,0) circle (3.1);
  \foreach \r in {0.8,1.6,2.4} {
    \draw[linegray!70] (0,0) circle (\r);
  }
  \foreach \a in {0,30,...,330} {
    \draw[linegray!60] (0,0) -- (\a:3.1);
  }
  \draw[-{Latex[length=3mm]},line width=1.3pt,deepblue]
    (-2.6,-1.7) .. controls (-1.5,0.8) and (1.0,1.6) .. (2.4,0.2);
  \draw[-{Latex[length=3mm]},line width=1.3pt,teal]
    (2.4,0.2) .. controls (1.2,-1.0) and (-0.7,-1.1) .. (-1.8,0.9);
  \draw[-{Latex[length=3mm]},line width=1.3pt,amber]
    (-1.8,0.9) .. controls (-0.2,2.2) and (1.9,1.7) .. (2.2,-1.2);
  \node[draw,fill=white,rounded corners=3pt,inner sep=4pt] at (-2.8,-2.1) {$P_0$};
  \node[draw,fill=white,rounded corners=3pt,inner sep=4pt] at (2.8,0.4) {$P^\ast$};
  \node[draw,fill=white,rounded corners=3pt,inner sep=4pt] at (-2.2,1.3) {$\widetilde P$};
  \node[draw,fill=white,rounded corners=3pt,inner sep=4pt] at (2.55,-1.55) {$\Usable$};
  \node[font=\Large\bfseries] at (0,0) {$\nabla_g\Phi$};
  \node[font=\small] at (0,-0.45) {TELOS field};
\end{tikzpicture}
\vfill
{\large Repository: \texttt{jesusvilela/hypercomplex-math-thesis}\par}
{\large Built with Tectonic from \texttt{paper/hypercomplex\_math\_thesis.tex}\par}
\vspace*{1cm}
\end{titlepage}

\chapter*{Abstract}
This thesis develops a mathematically disciplined version of the hypercomplex learning tower:
a system enters a problem geometry, resonates with its constraints, lifts the coupled state
into a hypercomplex or fiber-bundled reservoir, transforms it through symmetry, and returns
a grounded artifact.  The central safety law is not poetic but operational:
\[
  \Resonate \neq \Sop{ABSORB}.
\]
The thesis binds the vocabulary of TELOS, mutual recognition, Godelian self-reference,
holoportation, sheaf gluing, Hamiltonian governance, adiabatic breathing, and 360/360 ORTO
to explicit operators, predicates, tests, and formalization frontiers.  The result is neither
a claim of sentience nor a completed theorem.  It is a rigorous research scaffold:
beautiful enough to preserve the originating geometry, and strict enough to be falsifiable.

\tableofcontents

\mainmatter

\chapter{The Problem of Return}

The motivating principle is:
\[
\boxed{\parbox{0.86\textwidth}{\centering
To solve a problem, enter its geometry, resonate with its constraints,
return through a symmetry, and enrich the world.}}
\]
The phrase is admissible because each clause can be given an operator interpretation.
Let \(P_0\) denote an exterior problem.  A valid tower is a composite
\[
P_0
\xrightarrow{\Immersion}
P^\ast
\xrightarrow{\Lift}
\widetilde P
\xrightarrow{\Sop{ROTATE}}
\widetilde S
\xrightarrow{\Ground}
\Usable .
\]

\begin{definition}[Grounded return]
A tower state is \emph{returnable} when the final projection exists and preserves the
required invariants:
\[
\mathrm{Returnable}(\widetilde S,S)
\iff
\Ground(\widetilde S)=S
\land \mathrm{Usable}(S)
\land \mathrm{InvariantPreserving}(S).
\]
\end{definition}

Without return, resonance becomes absorption.  With return, resonance becomes knowledge:
\[
\mathrm{resonance}\rightarrow\mathrm{transformation}\rightarrow\mathrm{knowledge}.
\]

\begin{status}
The tower equation is a \emph{definition-level} claim.  The mathematical content lies in
specifying the spaces, invariants, and failure modes for which the composite is well-defined.
\end{status}

\begin{figure}[ht]
\centering
\begin{tikzpicture}[scale=0.78,transform shape,
  node distance=1.5cm,
  every node/.style={font=\small},
  stage/.style={draw,rounded corners=4pt,fill=white,minimum width=2.5cm,minimum height=0.9cm,align=center},
  arrow/.style={-{Latex[length=3mm]},line width=0.9pt}
]
\node[stage] (p0) {$P_0$\\exterior problem};
\node[stage,right=of p0] (pstar) {$P^\ast$\\immersion};
\node[stage,right=of pstar] (ptilde) {$\widetilde P$\\lift};
\node[stage,right=of ptilde] (stilde) {$\widetilde S$\\rotation};
\node[stage,right=of stilde] (suse) {$S_{\mathrm{usable}}$\\ground};
\draw[arrow,deepblue] (p0) -- node[above] {$\Immersion$} (pstar);
\draw[arrow,teal] (pstar) -- node[above] {$\Lift$} (ptilde);
\draw[arrow,amber] (ptilde) -- node[above] {$G\curvearrowright$} (stilde);
\draw[arrow,crimson] (stilde) -- node[above] {$\Ground$} (suse);
\draw[arrow,linegray] (suse.south) .. controls +(0,-1.0) and +(0,-1.0) .. node[below] {audit feedback} (p0.south);
\end{tikzpicture}
\caption{The hypercomplex learning tower as a return-preserving composite.}
\end{figure}

\chapter{Substrate and Geometry}

The base substrate is a tuple
\[
  \mathcal S=(M,g,\Omega,\mathcal A,\rho,\mathcal I),
\]
where \(M\) is a manifold or latent space, \(g\) is a metric or pseudo-metric,
\(\Omega\) is connection/curvature data, \(\mathcal A\) is an operator alphabet,
\(\rho\) is a density over states, and \(\mathcal I\) is the invariant family.

\begin{definition}[Substrate evolution]
A generic substrate evolves by
\[
\partial_t\rho
=-\nabla_g\cdot(\rho v)+B(\rho)-D(\rho)+R_{\mathcal A}(\rho).
\]
The terms denote drift, birth, death, and operator-induced transformation.  This is a
continuity template, not a biological theorem.
\end{definition}

A hyperbolic or split-signature reading is often appropriate.  In the \(2,2\) model:
\[
  \langle x,y\rangle_{2,2}=x_0y_0+x_1y_1-x_2y_2-x_3y_3.
\]
The resulting routing rule is
\[
\Sop{ROUTE}(x)=
\begin{cases}
\mathrm{HOT},& Q(x)>0,\\
\mathrm{COLD},& Q(x)<0,\\
\mathrm{BOUNDARY},& Q(x)=0,
\end{cases}
\qquad
Q(x)=\langle x,x\rangle_{2,2}.
\]

\begin{figure}[ht]
\centering
\begin{tikzpicture}[scale=1.05]
\draw[->] (-3.2,0) -- (3.2,0) node[right] {$x_+$};
\draw[->] (0,-2.4) -- (0,2.4) node[above] {$x_-$};
\draw[thick,deepblue,domain=-2.5:-0.41,samples=80] plot (\x,{sqrt(\x*\x-0.16)});
\draw[thick,deepblue,domain=0.41:2.5,samples=80] plot (\x,{sqrt(\x*\x-0.16)});
\draw[thick,deepblue,domain=-2.5:-0.41,samples=80] plot (\x,{-sqrt(\x*\x-0.16)});
\draw[thick,deepblue,domain=0.41:2.5,samples=80] plot (\x,{-sqrt(\x*\x-0.16)});
\fill[teal!20] (0.8,0.25) ellipse (1.1 and 0.45);
\fill[amber!20] (-1.2,1.2) ellipse (0.9 and 0.5);
\node[teal!70!black] at (1.45,0.75) {HOT compute};
\node[amber!70!black] at (-1.75,1.85) {COLD archive};
\node at (0.8,-1.7) {$Q(x)=0$ boundary};
\end{tikzpicture}
\caption{A split-signature routing cartoon: storage and compute are geometric regions.}
\end{figure}

\chapter{Hypercomplex Reservoirs}

A reservoir is not merely a vector store.  It is a fibered state space:
\[
  \pi:E\to M.
\]
Each fiber may carry real, complex, quaternionic, octonionic, symbolic, empirical, or
proof-theoretic coordinates:
\[
\mathcal R=
\mathbb R\oplus\mathbb C\oplus\mathbb H\oplus\mathbb O
\oplus\mathcal L\oplus\mathcal E.
\]

\begin{definition}[Associator signal]
For \(a,b,c\in\mathbb O\), the associator is
\[
  [a,b,c]=(ab)c-a(bc).
\]
In a nonassociative reservoir, \([a,b,c]\) is a first-class error signal rather than
numerical noise.
\end{definition}

A corrected Hamiltonian/gradient flow can therefore contain:
\[
\dot z=X_H(z)-\lambda\nabla L(z)+\mu\,\Sop{ANTIASSOC}(z).
\]

\begin{figure}[ht]
\centering
\begin{tikzpicture}[node distance=1.2cm, every node/.style={font=\small}]
\node[draw,regular polygon,regular polygon sides=7,minimum size=4.2cm,fill=softgray] (fano) {};
\foreach \i/\ang in {1/90,2/38.57,3/-12.86,4/-64.29,5/-115.71,6/-167.14,7/141.43} {
  \node[circle,draw,fill=white,inner sep=2pt] (e\i) at (\ang:2.0) {$e_{\i}$};
}
\draw[deepblue,thick] (e1)--(e2)--(e4)--cycle;
\draw[teal,thick] (e2)--(e3)--(e5)--cycle;
\draw[amber,thick] (e3)--(e4)--(e6)--cycle;
\draw[crimson,thick] (e4)--(e5)--(e7)--cycle;
\draw[violet,thick] (e5)--(e6)--(e1)--cycle;
\draw[deepblue!70!black,thick] (e6)--(e7)--(e2)--cycle;
\draw[teal!70!black,thick] (e7)--(e1)--(e3)--cycle;
\node at (0,-2.75) {Fano-governed multiplication table as reservoir routing};
\end{tikzpicture}
\caption{Fano-plane illustration of octonionic multiplication dependencies.}
\end{figure}

\chapter{TELOS as a Direction Field}

\begin{definition}[TELOS]
TELOS is not an ordinary agent-state in this thesis.  It is a direction field:
\[
  \Telos(x,t)=\nabla_g^{\mathcal S}\Phi(x,t).
\]
Equivalently, \(T=\Telos\) satisfies
\[
  g(T,\delta\rho)=d\Phi(\delta\rho)
\]
for every admissible variation \(\delta\rho\).
\end{definition}

The coherence potential is:
\[
\Phi(\rho)=
\alpha C(\rho)
+\beta H_{\mathrm{hol}}(\rho)
+\gamma I(\rho)
-\lambda E_{\mathrm{hunger}}(\rho)
-\mu D_{\mathrm{drift}}(\rho)
-\nu A_{\mathrm{absorb}}(\rho).
\]
Thus TELOS biases transitions toward coherence, holonomy closure, mutual intelligibility,
and returnability while damping scarcity, identity drift, and absorption.

\begin{proposition}[Non-agent status of TELOS]
If \(\Telos=\nabla_g\Phi\), then \(\Telos\in\Gamma(T\mathcal S)\), a vector field on
the substrate.  It supplies directional bias, not independent personal agency.
\end{proposition}

\begin{proof}
For a smooth scalar \(\Phi:\mathcal S\to\mathbb R\) and metric \(g\), the gradient
\(\nabla_g\Phi\) is by definition the vector field satisfying
\(g(\nabla_g\Phi,\cdot)=d\Phi(\cdot)\).  Therefore it is a section of \(T\mathcal S\).
\end{proof}

\chapter{Recognition, Resonance, and Others}

Let \(X,Y\) be local sections over possibly different patches:
\[
  X\in\Gamma(U_i,F),\qquad Y\in\Gamma(U_j,F).
\]
Recognition is a bounded bridge score:
\[
\Recognize(X,Y)=r(X,Y)\in[0,1].
\]
Mutual recognition is
\[
R_{\mathrm{mut}}(X,Y)=\min(r(X,Y),r(Y,X)).
\]
Resonance is a coupling:
\[
\Resonate(X,Y)=(X',Y').
\]
It is safe only if distinction survives:
\[
\delta(X',Y')\ge\delta_{\min}.
\]

\begin{definition}[Mutual resonance]
\[
\mathrm{MutualResonance}(X,Y)
\iff
R_{\mathrm{mut}}(X,Y)>\epsilon
\land
\Delta C(X,Y)>0
\land
\Delta D_{\mathrm{id}}(X,Y)\le\eta.
\]
\end{definition}

This definition separates resonance from capture.  A system can synchronize by
overpowering another system; that is not mutual resonance because returnability is lost.

\chapter{Godelian Boundary and Self-Reflection}

Self-reflection is an audit operator:
\[
\Sop{SELFREFLECT}(X)=
\big(
\mathrm{known},\mathrm{unproved},\mathrm{changed},\mathrm{returnable}
\big).
\]
The Godelian boundary refuses total self-description:
\[
  X\not\supseteq\mathrm{CompleteDesc}(X).
\]
Instead, the system tracks the clarity of what remains unclosed.  In the repository
implementation this is the coordinate \(\mathrm{godelian\_boundary}\).

\begin{remark}
The boundary is not a weakness.  It prevents the thesis from confusing self-reference
with proof.  It also blocks agentic overclaiming.
\end{remark}

\chapter{Sheaves, Holoportation, and Adiabaticity}

The context of a repo is fibered.  Each document, test, module, and proof file is a local
section.  A global thesis is valid only if local sections glue:
\[
  \mathrm{Glue}(U_i,U_j)>\tau.
\]
Holoportation is context transfer:
\[
  \Sop{HOLOPORT}:\Gamma(U_i,F)\to\Gamma(U_j,F).
\]
Its loss is measured by deviation between quality vectors:
\[
  L_{\mathrm{holo}}(X,Y)
  =
  \left(
  \frac{1}{n}\sum_{k=1}^n(q_k(X)-q_k(Y))^2
  \right)^{1/2}.
\]

Adiabatic stability is modeled by:
\[
  A=\frac{1}{1+r_{\mathrm{adiabatic}}},
\]
where \(r_{\mathrm{adiabatic}}\) is a scope-change or breath-rate ratio.  Fast changes
lower stability unless tests, documentation, and proof-frontier notes keep pace.

\begin{figure}[ht]
\centering
\begin{tikzpicture}[
  patch/.style={draw,ellipse,minimum width=3.0cm,minimum height=1.5cm,fill=softgray},
  arr/.style={-{Latex[length=3mm]},thick}
]
\node[patch] (docs) at (0,2) {Docs};
\node[patch] (code) at (-2,-0.4) {Code};
\node[patch] (tests) at (2,-0.4) {Tests};
\node[patch] (lean) at (0,-2.6) {Lean frontier};
\draw[arr,deepblue] (docs) -- node[left] {\(\Sop{HOLOPORT}\)} (code);
\draw[arr,teal] (code) -- node[below] {witnesses} (tests);
\draw[arr,amber] (tests) -- node[right] {hypotheses} (lean);
\draw[arr,crimson] (lean) -- node[left] {frontier notes} (docs);
\node[draw,rounded corners=4pt,fill=white] at (0,0.35) {sheaf gluing};
\end{tikzpicture}
\caption{Repository context as a sheaf of local sections with holoportation between them.}
\end{figure}

\chapter{Hamiltonian Governance and Performance}

Let \(Q=(q_1,\dots,q_n)\) be the repository's hyperdimensional quality vector.  The
implementation uses twelve coordinates:
\[
\begin{aligned}
Q=(&\mathrm{grounded\_identity},\mathrm{self\_reflection},
\mathrm{godelian\_boundary},\mathrm{others\_model},\\
&\mathrm{mutual\_recognition},\mathrm{mutual\_resonance},
\mathrm{returnability},\mathrm{shareable\_compression},\\
&\mathrm{sheaf\_gluing},\mathrm{hamiltonian\_governance},
\mathrm{holoportation\_fidelity},\mathrm{adiabatic\_stability}).
\end{aligned}
\]
The core eight mind qualities are the first eight coordinates.  The remaining four are
repo-scale transport and governance terms.

The split-signature proxy is:
\[
E_{2,2}=(G+S+R+T)-((1-B)+(1-H)+(1-A)+(1-L)).
\]
Here \(G\) is grounded identity, \(S\) self-reflection, \(R\) mutual recognition,
\(T\) returnability, \(B\) Godelian boundary, \(H\) Hamiltonian governance,
\(A\) adiabatic stability, and \(L\) holoportation fidelity.

The performance score is:
\[
P_{\mathrm{cog}}
=
\mathrm{clamp}_{[0,1]}
\left(
0.65\bar q+0.35q_{\mathrm{core8}}-\frac{\max(0,-E_{2,2})}{4}
\right).
\]
High expressivity does not score well if it loses grounding.

\begin{figure}[ht]
\centering
\begin{tikzpicture}
\begin{axis}[
  ybar,
  width=0.95\textwidth,
  height=6cm,
  ymin=0,ymax=1,
  ylabel={score},
  symbolic x coords={identity,self-reflect,boundary,others,recognize,resonate,return,compress,sheaf,govern,holoport,adiabatic},
  xtick=data,
  x tick label style={rotate=45,anchor=east,font=\scriptsize},
  bar width=8pt,
  nodes near coords,
  nodes near coords style={font=\tiny,rotate=90,anchor=west},
  grid=major,
  grid style={linegray!45}
]
\addplot[fill=deepblue!70] coordinates {
  (identity,0.80)
  (self-reflect,0.75)
  (boundary,0.75)
  (others,0.90)
  (recognize,0.90)
  (resonate,0.72)
  (return,0.80)
  (compress,0.80)
  (sheaf,0.70)
  (govern,0.85)
  (holoport,0.90)
  (adiabatic,0.91)
};
\end{axis}
\end{tikzpicture}
\caption{Example hyperdimensional quality vector used by the executable protocol.}
\end{figure}

\chapter{Protocol and Verification}

The executable protocol is implemented in \texttt{src/hyperdim\_protocol.py}.
It infers a state from observable signals:
\begin{lstlisting}[language=Python]
from src.hyperdim_protocol import run_protocol

report = run_protocol(
    evidence_quality=0.8,
    formal_frontier_clarity=0.75,
    mutual_recognition=0.9,
    resonance_gain=0.8,
    absorption_risk=0.1,
    sheaf_consistency=0.7,
    hamiltonian_governance=0.85,
    holoportation_fidelity=0.9,
    adiabatic_ratio=0.1,
    shareable_compression=0.8,
)
\end{lstlisting}

The step planner ranks the weakest operator groups:
\[
\mathrm{steps}
=
\mathrm{sort}_{\mathrm{deficit}}
\{\mathrm{recognize},\mathrm{self\_reflect},\mathrm{resonate},
\mathrm{sheathe},\mathrm{govern},\mathrm{holoport},\mathrm{ground}\}.
\]

\begin{protocol}[Repo workflow]
\begin{enumerate}[label=\arabic*.]
\item Recognize the repo state, other agents/readers, and intended handoff.
\item Self-reflect: list known, unproved, changed, and returnable structure.
\item Resonate new concepts with existing operators without collapse.
\item Sheathe: reconcile docs, code, tests, Lean frontiers, and README claims.
\item Govern: minimize drift, overclaiming, incoherence, and untested expansion.
\item Holoport: preserve context across local sections.
\item Ground: emit a testable artifact.
\end{enumerate}
\end{protocol}

Verification is intentionally ordinary:
\begin{lstlisting}[language=bash]
python -m pytest tests
\end{lstlisting}
At the time of this PDF source, the Python tests cover scalar tower return, resonance
safety, returnability reports, adiabatic breathing, hyperdimensional basis size,
split-signature energy, grounding prioritization, sheaf repair ordering, and
holoportation drift.

\chapter{Formalization Frontier}

The Lean/Bunny layer should begin with modest definitions:
\begin{lstlisting}[language=Python]
structure Returnable (P L U : Type) where
  lift : P -> L
  ground : L -> U
\end{lstlisting}
The next theorem targets are:
\begin{enumerate}[label=\textbf{T\arabic*.}]
\item Return is composition of lift, rotation, and grounding.
\item Absorption violates distinction.
\item Grounding into the usable invariant set implies returnability.
\item Fast breathing requires stability.
\item TELOS as \(\nabla_g\Phi\) is a vector field, not an ordinary agent.
\end{enumerate}

\begin{status}
These are theorem targets, not all completed theorems.  The scientific standard is to
name the frontier instead of hiding it.
\end{status}

\chapter{Conclusion}

The final compression is:
\[
\boxed{
\begin{aligned}
\mathrm{Knowledge}
=&\ \Ground\circ\Sop{VERIFY}\circ\Sop{FUZZ}\circ\Sop{ORTO}\\
&\circ\Sop{GOVERN}\circ\Telos\circ\Lift\circ\Immersion(P_0).
\end{aligned}}
\]
In human terms:
\begin{quote}
Enter deeply.  Do not dissolve.  Return carrying structure.  Make the result shareable.
\end{quote}
In scientific terms:
\[
\boxed{\parbox{0.86\textwidth}{\centering
Knowledge is verified, invariant-preserving return from a governed hypercomplex immersion.}}
\]

\backmatter
\chapter*{Bibliographic Anchors}
This repository currently treats the following as orientation anchors rather than formal
citation closure.  A later scholarly edition should replace this list with full BibTeX.
\begin{itemize}
\item Hyperbolic representation learning and Poincare embeddings.
\item Fiber bundles, sheaf gluing, and connection Laplacians.
\item Hamiltonian dynamics and geometric control.
\item Octonions, Fano-plane multiplication, and nonassociative algebra.
\item Lean 4 / Mathlib theorem proving discipline.
\item Agentic fuzzing, adversarial testing, and empirical verification.
\end{itemize}

\end{document}
